\chapter*{ДҮГНЭЛТ}
\addcontentsline{toc}{chapter}{ДҮГНЭЛТ}
Энэхүү бакалаврын судалгааны ажлын хүрээнд Sumbee.mn сошиал мониторингийн системд тулгарч буй хэрэгцээ шаардлага болох агуулгын шинжилгээний нарийвчлалыг сайжруулах зорилгоор Topic Modeling болон NER (Named Entity Recognition)-ийн нэгтгэсэн шийдлийг боловсрууллаа. 

Боловсруулсан загварыг инференс пайплайнаар ажиллуулж туршсан үр дүн нь энгийн мэдлэгийн сан (Knowledge Database) үүсгэхээс илүү контекст таних чадвар сайтай, LLM API ашигласнаас хэд дахин хямд өртөгтэй бөгөөд хурдтай гэдгийг батлав. Уг шийдэл нь сошиал тренд болон үүсч буй чухал сэдвүүдийг автоматжуулсан байдлаар гаргахад чухал ач холбогдолтой юм.

Цаашид олон хэлний загвар (Multilingual) болон монгол хэл дээр тусгайлан fine-tune хийгдсэн Topic Classification загварыг аппликэйшнд бодит хугацааны инференст оруулах, сургагдсан моделуудыг API хэлбэрт оруулан Sumbee.mn-ийн микросервис архитектурт бүрэн нэгтгэх боломжтой юм. Мөн Network graph-д temporal (цаг хугацааны) анализ нэмж, чиг хандлагын өөрчлөлтийг урьдчилан таамагладаг болгохоор төлөвлөгдлөө.
